\subsection{Independent re-proof and the configuration proof comparison}

The independent proof was written from the statement of A.CF2 and the
hypothesis of A.CF1 before either earlier proof or any earlier audit was
opened. Its preserved source is \path{work/configuration/blind_proof.tex},
with SHA-256
\begin{center}\footnotesize\ttfamily
d435494f67de82bd06f08a04cc519501e739a7abf0f88dc3b8b42ef5c96c6ac9.
\end{center}
The comparison then checked every line 1--140 of
\path{background/A1_checkpoint03/sections/A_configuration.tex}; the
line-level record is \path{work/configuration/S1_line_comparison.tsv}.
This is an independent re-proof by GPT-6 Astra Pro, followed by a
targeted line-by-line comparison. It is not a claim of review by a
different model family. The run is in Pro mode; a separate effort
setting is not exposed.

The proofs agree. The independent proof detects finite fibre quotients
inside finite automorphism groups rather than first passing to the
whole profinite automorphism group. To make this version completely
self-contained, we record its group-theoretic details.
Let $P$ be a nonabelian finitely generated free pro-$r$ group, and let
$U_n$ be a decreasing characteristic open basis. Put $P_n=P/U_n$ and
\[
 Z_n=\{x\in P:[x,P]\subseteq U_n\}.
\]
The finite group
$\ker(\operatorname{Aut}(P_n)\to\operatorname{Aut}(P_n/\Phi(P_n)))$
is an $r$-group: it acts freely on the ordered lifts of a basis of
$P_n/\Phi(P_n)$, a set of cardinality $|\Phi(P_n)|^d$.
Each lift generates $P_n$, since otherwise its generated subgroup is
contained in a maximal subgroup and its image cannot span the
Frattini quotient. Freeness of the action and the orbit formula give
the claim. Thus a group acting through an $r$-group on $P/\Phi(P)$
has $r$-group image in every $\operatorname{Aut}(P_n)$.

The centre of $P$ is trivial. For completeness the Magnus homomorphism
sends a free basis $x_i$ to $1+X_i$ in
$1+(X_1,\ldots,X_d)\mathbf F_r\langle\!\langle X_1,\ldots,X_d\rangle\!\rangle$.
Its target is pro-$r$ under degree truncation, so the map exists by
the universal property of $P$. For every finite $r$-quotient $Q$,
substitution $X_i\mapsto q_i-1$ recovers the quotient in
$\mathbf F_r[Q]^\times$. The augmentation ideal is nilpotent:
choose a central element $z$ of order $r$, apply induction to
$Q/\langle z\rangle$, and use $(z-1)^r=0$; if the quotient
augmentation ideal has exponent $m$, the original ideal has exponent
at most $mr$. Hence the substitution is defined on formal series,
and separation by finite quotients proves injectivity of Magnus.
A homogeneous polynomial commuting with $X_1$ has only the monomial
$X_1^n$: a monomial beginning in another letter would contribute to
$fX_1$ but not to $X_1f$; factor out $X_1$ and repeat.
Commutation with $X_2$ eliminates every term of positive degree.
A central Magnus unit therefore equals $1$.

Consequently $\bigcap_nZ_n=Z(P)=1$. Given an open normal $V\subset P$,
compactness implies $Z_n\subset V$ for some $n$. The finite
$r$-group image of $H$ in $\operatorname{Aut}(P_n)$ has kernel on
the fibre contained in $V$. This proves A.CF1 and its domination
assertion without assuming residual-$r$ properties of the base.
It also explains why using one prescribed finite fibre quotient
alone would be insufficient: its centre could be lost. The earlier
proof avoids that loss by using the whole centreless $P$, and is
equally valid. The earlier proof supplies a more explicit argument
for Magnus injectivity than the initially saved blind proof; the
argument just given incorporates and verifies that detail.

For A.CF2, the independent proof wrote the peripheral homology module
explicitly as
\[
 M_A=\mathbf F_r[A]/\mathbf F_r\Bigl(\sum_{a\in A}e_a\Bigr).
\]
The vector in the denominator is nonzero in every characteristic.
The monodromy is trivial on $M_A$ and on $H_1(C',\mathbf F_r)$,
and hence consists of shears. This verifies the earlier assertion
even when $r\mid|A|$; no averaging or equivariant splitting appears
in either proof. The earlier choice $A=Q/R$ is sufficient; the blind
proof permits any larger proper quotient, which gives the same
conclusion and can dominate several fibre quotients simultaneously.

Both proofs use the normal cover with deck group $A\times A$, the
identity $\Lambda\cap\pi_1(U_A)=\pi_1(C'-a^{-1}(P))$, and the finite
core over the cosets of that normal subgroup. The normal core has
$r$-group kernel over $A\times A$, because it embeds in a product of
conjugate $r$-group quotients. The image of $K_E$ lies in this kernel,
so the quotient factors through $E_{\mathrm{cr}}$. The simultaneous
finite image of the maps to $T$ and $Q$ verifies that fibre domination
persists after completion. Finally both proofs derive middle and
right exactness from the finite-quotient universal property of
pro-$\Sigma$ completion. There is no appeal to left exactness of a
completion functor. We found no error in A.CF1 or A.CF2.

\paragraph{Check record.}
A.CF1 and A.CF2 were also checked in this writing, by independent
re-proof before access to the earlier proofs and by subsequent
line-by-line comparison, on 18 November 2026. Their earlier checks
are retained in the inherited check ledger. The direct
characteristic-$p$ argument A.DC1--A.DC4 was checked only in this
writing; it is logically separate from this blind comparison.
